\subsection*{Общая характеристика работы}

\newcommand{\actuality}{\underline{\textbf{Актуальность темы.}}}
\newcommand{\aim}{\underline{\textbf{Целью}}}
\newcommand{\tasks}{\underline{\textbf{задачи}}}
\newcommand{\defpositions}{\underline{\textbf{Основные положения, выносимые на~защиту:}}}
\newcommand{\novelty}{\underline{\textbf{Научная новизна:}}}
\newcommand{\influence}{\underline{\textbf{Практическая значимость}}}
\newcommand{\reliability}{\underline{\textbf{Достоверность}}}
\newcommand{\probation}{\underline{\textbf{Апробация работы.}}}
\newcommand{\contribution}{\underline{\textbf{Личный вклад.}}}
\newcommand{\publications}{\underline{\textbf{Публикации.}}}

{\actuality}
В настоящее время в Институте ядерной физики им. Г. И. Будкера СО РАН проводятся эксперименты на электрон-позитронном накопителе \text{ВЭПП-2000} \cite{VEPP} в диапозоне энергий от 320 МэВ до 2 ГэВ в системе центра масс. В местах встречи пучков установлены два детектора: криогенный магнитный детектор \text{КМД-3 \cite{CMD}} и сферический нейтральный детектор \text{СНД \cite{SND}.} Физическая программа экспериментов включает в себя измерение сечений электрон-позитронной аннигиляции в адроны на низких энергиях, а также измерение адронного вклада в вычисление аномального магнитного момента мюона {$(g-2)_\mu$}. Так как в данной области энергии расчеты в рамках пертурбативной квантовой хромодинамики невозможны, то основным источником изучения взаимодействия легких кварков являются данные, полученные в экспериментах. Набор экспериментальных данных с детектором \text{КМД-3} начат в декабре 2010 года. В апреле 2013 года эксперименты были остановлены в связи с модернизацией инжекционного комплекса \cite{status}. В конце 2016 года запущен новый инжекционный комплекс \text{ВЭПП-5 \cite{VEPP-5}} и возобновлен набор статистики детектором \text{КМД-3.}

 \aim\ данной работы является измерение сечения процесса {$e^+\:e^- \to K_{S}\:K_{L}\:\pi^0$}
 в области энергий от порога рождения \text{(1130 МэВ)} до \text{2 ГэВ.} В анализе процесса использовались данные, полученные в экспериментах 2011-2012 годов с детектором КМД-3. Интеграл светимости составил \text{$33.18$ пб$^{-1}$} в диапозоне энергий от \text{1.1 ГэВ} до \text{2 ГэВ. }
 %Процесс примечателен тем, что все три конечные частицы псевдоскалярные и имеют общую плоскость распада. Одно из главных физических свойства процесса - это нейтральность конечных частиц. (Ниже воткнуть.)
 %пересчитать точнее светимость для твоей области

В задачи данной работы входило:
\begin{enumerate}
  \item Выработать оптимальные критерии отбора событий для процесса {$e^+\:e^- \to K_{S}\:K_{L}\:\pi^0$}.
  \item Определить эффективность регистрации процесса с помощью Монте-Карло моделирования.
  \item Получить предварительное сечение процесса {$e^+\:e^- \to K_{S}\:K_{L}\:\pi^0$} в диапозоне энергий от порога рождения \text{(1130 МэВ)} до \text{2 ГэВ.}
\end{enumerate}

%\novelty\ данной работы заключается в исследовании малоизученного процесса  {e^+\:e^- \to $\itshape K_{S}\:K_{L}\:\pi^0$}, а так же измерение его вклада в поляризацию ваакума, используемое для вычисления аномального магнитного момента мюона {$(g-2)_\mu$}.

Сечение данного процесса до недавнего времени не было измерено, в настоящее время существуют пока что только предварительные результаты, представленные коллаборацией SND \cite{korneev}, и недавно опубликованная статья коллаборации BaBar \cite{BaBar}. %Результаты ранее сказанных коллабораций имеют различия в области энергий выше 1.8 ГэВ, таким образом, имеет место дополнительная задача --- подтвердить результаты одной из коллабораций.
%данной работы заключается в вычислении сечения процесса  {e^+\:e^- \to $\itshape K_{S}\:K_{L}\:\pi^0$} в настоящее время предварительно изммеренного коллаборациями BaBar и SND, а так же в измерении его вклада в поляризацию ваакума, используемое для вычисления аномального магнитного момента мюона {$(g-2)_\mu$}. %коряво

%Научная новизна заключается в том, что сечение данного процесса до сих пор не был измерен, в настоящее время существуют пока что только предварительные результаты, представленные коллаборациями Бабар и СНД.
%Сказать про физику здесь, изотопический спин, все частицы в одной плоскости, три псевдоскалярные частицы. % Характеристика работы по структуре во введении и в автореферате не отличается (ГОСТ Р 7.0.11, пункты 5.3.1 и 9.2.1), потому её загружаем из одного и того же внешнего файла, предварительно задав форму выделения некоторым параметрам

%Диссертационная работа была выполнена при поддержке грантов ...

%\underline{\textbf{Объем и структура работы.}} Диссертация состоит из~введения, четырех глав, заключения и~приложения. Полный объем диссертации \textbf{ХХХ}~страниц текста с~\textbf{ХХ}~рисунками и~5~таблицами. Список литературы содержит \textbf{ХХX}~наименование.

%\newpage
\subsection*{Содержание работы}
Во \underline{\textbf{введении}} обосновывается актуальность исследований, проводимых в рамках данной диссертационной работы, приводится обзор научной литературы по изучаемой проблеме, формулируется цель, ставятся задачи работы, сформулированы научная новизна и практическая значимость представляемой работы.

\underline{\textbf{Первая глава}} посвящена ...

 картинку можно добавить так:
\begin{figure}[ht] 
  \center
  \includegraphics [scale=0.27] {latex}
  \caption{Подпись к картинке.} 
  \label{img:latex}
\end{figure}

Формулы в строку без номера добавляются так:
\[ 
  \lambda_{T_s} = K_x\frac{d{x}}{d{T_s}}, \qquad
  \lambda_{q_s} = K_x\frac{d{x}}{d{q_s}},
\]

\underline{\textbf{Вторая глава}} посвящена исследованию 

\underline{\textbf{Третья глава}} посвящена исследованию 

В \underline{\textbf{четвертой главе}} приведено описание 

В \underline{\textbf{заключении}} приведены основные результаты работы, которые заключаются в следующем:
%% Согласно ГОСТ Р 7.0.11-2011:
%% 5.3.3 В заключении диссертации излагают итоги выполненного исследования, рекомендации, перспективы дальнейшей разработки темы.
%% 9.2.3 В заключении автореферата диссертации излагают итоги данного исследования, рекомендации и перспективы дальнейшей разработки темы.
\begin{enumerate}
  \item Были выработаны оптимальные критерии отбора событий процесса ${e^+\:e^- \to K_{S}\:K_{L}\:\pi^0}$. Вклад данного процесса в общее адронное сечение менее 5\%, тем не менее удалось набрать статистику в 1924 событий при помощи выбора оптимальных условий отбора.
  \item Определена эффективность регистрации при помощи Монте-Карло моделирования, которая составила около 2\% в диапозоне энергий от \text{1.1 ГэВ} до \text{1.4 ГэВ} и около 4.5\% для энергий от 1.4 ГэВ до 2 ГэВ со статистической ошибкой $\sim1\%$.
  \item Измерено сечение процесса ${e^+\:e^- \to K_{S}\:K_{L}\:\pi^0}$ в диапозоне энергий от \text{1.1 ГэВ} до \text{2 ГэВ} со статистической погрешностью около 10\%. Вычислены радиационные поправки с точностью $\sim1\%$, с последующим получением конечного результата --- борновское сечение изучаемого процесса.
  \item Сделана оценка вклада трех основных фоновых событий, общий вклад которых составил не более 10\% от сечения процесса ${e^+\:e^- \to K_{S}\:K_{L}\:\pi^0}$.
\end{enumerate}



%\newpage
\renewcommand{\refname}{\large Публикации автора по теме диссертации}
%\nocite{*}
%\insertbiblioauthor                          % Подключаем Bib-базы
\insertbibliofull
